% === Section 1: Introduction ===
\section{Introduction}\label{sec:intro}

The deployment of AI-driven controllers in safety-critical physical systems---battery
energy storage, autonomous vehicles, surgical robotics, power grid dispatch---has
progressed from laboratory prototypes to field-deployed assets.  A pervasive but
under-addressed hazard accompanies this transition: the controller takes actions
conditioned on \emph{observations} of the true physical state, yet observations are
never perfect.  Sensors stall, telemetry packets arrive late or out-of-order,
calibration drifts silently, and network links drop entire channels.  When the
observation pipeline degrades, the gap between what the controller \emph{believes} to
be the state and what the state \emph{actually is} may grow unboundedly, and the
controller may take an action that is formally safe in the observed state space yet
catastrophically unsafe in the true state space.

We call this gap the \textbf{Observation--Action Safety Gap} (OASG).  The OASG is
distinct from standard notions of model uncertainty, adversarial perturbation, or
out-of-distribution detection.\footnote{Out-of-distribution (OOD) detectors identify
distributional mismatches but do not synthesise conservative action sets; adversarial
robustness certifies bounded input perturbations but does not handle telemetry
infrastructure faults (packet loss, reordering, drift).}  It requires a runtime-level
mechanism that:
\begin{enumerate}[nosep]
  \item Detects telemetry degradation at every control step.
  \item Quantifies the resulting \emph{set of states consistent with the degraded
        observation}.
  \item Tightens the admissible action space so that safety is guaranteed for
        every member of that set.
  \item Shields or repairs any candidate action that exits the tightened set.
  \item Issues an auditable, tamper-evident certificate for every released action.
\end{enumerate}

No existing framework---simplex architectures~\cite{sha2001using},
runtime-verification monitors~\cite{leucker2009brief}, Bayesian OOD
detectors~\cite{hendryx2019deep}, or conformal-prediction wrappers~\cite{romano2019conformalized}---addresses
all five requirements simultaneously while remaining \emph{domain-agnostic},
\ie~applicable to any cyber-physical domain via a thin adapter rather than a domain-specific redesign.

\subsection{The ORIUS Framework}

This paper presents \textbf{ORIUS} (Observation--Reality Integrity for Universal
Safety), a typed runtime safety layer whose five-stage kernel---\textbf{Detect},
\textbf{Calibrate}, \textbf{Constrain}, \textbf{Shield}, \textbf{Certify}, abbreviated
\dcs{}---closes the OASG for any upstream controller.  The kernel is defined over abstract
typed contracts (\cref{sec:contracts}); domain specifics are confined to an adapter
module that implements six typed methods.  A formal theorem chain (T1--T11) connects
the runtime invariants to end-to-end safety guarantees under stated assumptions
(A1--A8).

\subsection{Scope and Claims}

Unlike prior conference publications that motivated ORIUS with six candidate domains,
this manuscript restricts its empirical claims to the \textbf{two domains fully
validated end-to-end} with locked, reproducible artifacts:

\begin{enumerate}[nosep]
  \item \textbf{Battery Energy Storage (BESS)}: dispatch optimisation under degraded
        state-of-charge telemetry, with DeepOQE neural forecasting and CPSBench 168-hour
        replay.
  \item \textbf{Autonomous Vehicles (AV)}: longitudinal prediction of ego-speed and
        relative gap under degraded perception, validated on the full Waymo Open Motion
        Dataset corpus (7~shards, 1{,}975~scenarios).
\end{enumerate}

Both domains are validated with the \emph{same kernel binary}, the same contract
stack, the same governance pipeline, and the same CertOS certificate chain.  The
remaining four candidate domains (industrial monitoring, healthcare, maritime
navigation, aerospace) are discussed only as future directions.

\subsection{Contributions}

The principal contributions of this paper are:

\begin{enumerate}[nosep]
  \item A complete formal architecture for runtime safety under degraded observation,
        expressed as five typed stages with a supporting contract algebra
        (\cref{sec:architecture}).
  \item An 11-theorem proof chain connecting observation reliability to true-state
        safety guarantees (\cref{sec:theory}).
  \item A shift-aware adaptive uncertainty mechanism that dynamically adjusts
        prediction intervals without offline recalibration (\cref{sec:shift}).
  \item Equal-depth empirical validation on two disparate CPS domains:
        \begin{itemize}[nosep]
          \item \textbf{Battery}: 288 runtime steps across 4 controller variants,
                zero true-state violations (TSVR = 0), 100\% audit completeness.
          \item \textbf{AV}: 9{,}348 runtime steps across 228 test scenarios,
                TSVR reduced from 0.660 (baseline) to 0.628 (ORIUS), intervention
                rate 90.2\%, 100\% audit completeness.
        \end{itemize}
  \item A fully reproducible artifact bundle with SHA-256 hashes, locked
        dependencies, and a single-command pipeline (\cref{sec:reproducibility}).
\end{enumerate}

\paragraph{What this paper proves and what it does not.}
The theorem chain (T1--T11) is proved in full generality over the typed contract
algebra; these results hold for \emph{any} domain that satisfies assumptions A1--A8.
Empirical confirmation, however, is restricted to the two domains above.  In
particular:
\begin{itemize}[nosep]
  \item The Battery domain achieves $\TSVR = 0$ (zero true-state violations) and
        serves as a \emph{reference witness}: it demonstrates the full pipeline
        end-to-end with the deepest artifact-to-theorem closure.
  \item The AV domain is a \emph{bounded proof-validated} row: ORIUS reduces the
        TSVR from 0.660 to 0.628 under a longitudinal TTC + predictive-entry-barrier
        contract.  A residual TSVR floor remains (see \cref{sec:limitations}).
  \item Four additional candidate domains (industrial monitoring, healthcare,
        maritime navigation, aerospace) motivate the architecture but have
        \textbf{no locked pipeline artifacts} in this release.  They appear only in
        future-work discussion and are not part of any empirical claim.
  \item Shift-awareness guarantees are bounded to the drift rates observed during
        validation; adversarial robustness under intentional telemetry attacks has
        not been evaluated.
\end{itemize}

\subsection{Paper Organisation}

\Cref{sec:problem} formalises the OASG problem.  \Cref{sec:related} surveys related
work.  \Cref{sec:architecture} details the \dcs{} kernel and typed interfaces.
\Cref{sec:contracts} covers the contract stack and governance surface.
\Cref{sec:theory} presents the theorem bridge connecting contracts to safety.
\Cref{sec:shift} introduces the shift-aware mechanism.  \Cref{sec:battery}
and~\cref{sec:av} provide comprehensive empirical validation on BESS and AV,
respectively.  \Cref{sec:crossdomain} synthesises cross-domain findings.
\Cref{sec:reproducibility} documents the reproducibility infrastructure.
\Cref{sec:limitations} discusses limitations.  \Cref{sec:conclusion} concludes.
